% BHU PYQ -- Professional Exam Template
\documentclass[11pt,a4paper]{article}

\usepackage[a4paper,top=2.5cm,bottom=2.5cm,left=2.8cm,right=2.8cm,headheight=14pt]{geometry}
\usepackage{amsmath,amssymb}
\usepackage[T1]{fontenc}
\usepackage[utf8]{inputenc}
\usepackage{lmodern}
\usepackage{microtype}
\usepackage{fancyhdr}
\usepackage{enumitem}
\usepackage{listings}
\usepackage{hyperref}
\usepackage{lastpage}
\usepackage{parskip}

\hypersetup{colorlinks=true,urlcolor=black,linkcolor=black,
  pdftitle={CS-103: Numerical Computing -- BHU PYQ}}

\pagestyle{fancy}
\fancyhf{}
\renewcommand{\headrulewidth}{0.4pt}
\renewcommand{\footrulewidth}{0.4pt}
\fancyhead[L]{\small\textit{Banaras Hindu University}}
\fancyhead[R]{\small\textit{CS-103: Numerical Computing}}
\fancyfoot[C]{\small Page \thepage\ of \pageref{LastPage}}

\lstset{
  basicstyle=\small\ttfamily,
  frame=single,
  breaklines=true,
  columns=flexible,
  keepspaces=true
}

\newlist{parts}{enumerate}{2}
\setlist[parts,1]{label=(\alph*),leftmargin=*,itemsep=4pt,topsep=3pt}
\setlist[parts,2]{label=(\roman*),leftmargin=*,itemsep=2pt,topsep=2pt}
\newcommand{\pts}[1]{\hfill\textit{[#1]}}

\begin{document}
\begin{center}
  {\large\bfseries BANARAS HINDU UNIVERSITY}\\[2pt]
  {\normalsize B.Sc. (Hons.) Semester III Examination, 2023-24}\\[6pt]
  \rule{0.8\linewidth}{0.5pt}\\[6pt]
  {\large\bfseries Paper: CS-103 --- Numerical Computing}\\[4pt]
  {\normalsize Time: Three hours \hfill Maximum Marks: 70}
\end{center}

\rule{\linewidth}{0.4pt}

\smallskip
\noindent\textit{Instructions: Attempt ANY FIVE questions.}

\bigskip

ROUGIN OS pocsssesn so sccteocsseebeveenas
CS-103 : Numerical Computing

\medskip\noindent\textbf{Question 1.} a. Diagonalize the given matrix: :\pts{7}
52; O22

\medskip\noindent\textbf{Question 0.} 3220
t \& -
\begin{parts}
  \item Solve the following equation using bisection method to find a real root between 2 and 3
  (correct to 3 decimal places)\pts{7}
  xlogigx —1.2=0 , .
\end{parts}

\medskip\noindent\textbf{Question 2.} For the following system of linear equations:
9x, + 2x2 +4x3 = 20 ,
X, + 10x. + 4x3 =6
2x, — 4x2 + 10x3 = -15
a: Check the convergence of the iterative procedure. [4]
se \_b. Use the Gauss-Siedel iterative procedure to solve the above (correct to 3 significant digits).\pts{10}

\medskip\noindent\textbf{Question 3.} a. Derive the Lagrange interpolation formula. What are the advantages and dis-advantages ofusingit? [7] \_
\begin{parts}
  \item Fit 3-degree Lagrange polynomial to find f(x) at x = 3.8 for the following function:\pts{7}
  ASHE Saad ce ER Se
  f(x 6 Chae TE 2 Saee
\end{parts}

\medskip\noindent\textbf{Question 4.} .a. Use the Newton’s divided difference formula to Gerive; :\pts{7}
i, Trapezoidal Rule
  \begin{parts}
    \item Simpson’s 1/3 Rule :
  \end{parts}
\begin{parts}
  \item Use Gauss-Legendre 2-point formula to calculate:\pts{7}
  1
  Je — 10x + 6 dx
\end{parts}

\medskip\noindent\textbf{Question 4.}

\medskip\noindent\textbf{Question 5.} Solve the following equations (correct to 4 significant digits) using:\pts{14}
d
2x, y(0)=1, 0<x<05, h=01
dx :
i, Euler’s Method
ii, Picard’s Method
Pe. Te\pts{0}

\medskip\noindent\textbf{Question 6.} For the given matrix calculate:\pts{14}
\begin{parts}
  \item Inverse matrix using Gauss Elimination method.
  \begin{parts}
    \item Largest Eigen value and corresponding Eigen vector using iterative method,
    1 <3: 4
    S152
  \end{parts}
\end{parts}

\medskip\noindent\textbf{Question 2.} 40
| . 7. Compute the root between | and 2 of the equation given below using:\pts{14}
x3 — 2x2 437-5 =0
\begin{parts}
  \item Regula-Falsi method :
  \begin{parts}
    \item Newton Raphson method
    OO
    | ee ee eee ey Se
  \end{parts}
\end{parts}

\vfill
\rule{\linewidth}{0.4pt}
\begin{center}\small
\href{https://bhu-exam.vercel.app/}{Click here to visit our website}\\[4pt]\href{https://me-aryan.vercel.app/}{Click here to visit our contributor}\\[4pt]\href{https://quashan-me.vercel.app/}{Click here to visit our contributor}
\end{center}
\end{document}
